\clearpage\begingroup\SourceLeaf{66b}\HistoricalHead{61}{R}{BERECHNUNG MAGNETISCHER KRÄFTE.}
\begin{center}
[3.]\par\medskip\rule{.86\textwidth}{.35pt}\\
{\small [Aus Handbuch 19, Be, Kleine Aufsätze aus verschiedenen Theilen der Mathematik.\\
Angefangen im März 1809, S. 185.]}\\
\rule{.86\textwidth}{.35pt}
\end{center}
\noindent\hspace*{\parindent}Die eine Bedingung zwischen Declinationen und Intensitäten kann am
einfachsten analytisch so dargestellt werden.

Es sein die Coordinaten von drei Orten \(x,y;\ x',y';\ x'',y''\); die partiellen
horizontalen magnetischen Kräfte an denselben \(\xi,\eta;\ \xi',\eta';\ \xi'',\eta''\). Dann ist
\[
\begin{aligned}
0={}&\xi x'+\xi'x''+\xi''x-\xi x''-\xi'x-\xi''x'\\
&{}+\eta y'+\eta'y''+\eta''y-\eta y''-\eta'y-\eta''y'.
\end{aligned}
\]
\noindent Geometrisch so:
\par\indent Es sein beim Durchschreiten des Umfangsdreiecks, welches die drei
Punkte 0, 1, 2 bilden, die Drei Seiten \(r'',r,r'\), deren Azimuthe \(A'',A,A'\); die
Drei Declinationen \(\delta,\delta',\delta''\); die Drei horizontalen Intensitäten \(h,h',h''\), so ist
\[
0=rh\cos(A-\delta)+r'h'\cos(A'-\delta')+r''h''\cos(A''-\delta'').
\]
\noindent NB. die Declinationen sind gegen Einen bestimmten Meridian zu nehmen.
\bigskip\begin{center}\textbf{BEMERKUNGEN.}\end{center}
{\small\noindent\hspace*{\parindent}Wird die Horizontalkomponente der magnetischen Kraft ihrer Größe und Richtung nach durch die
Vektorbezeichnung \(\mathfrak H\) charakterisiert, das Linienelement einer beliebigen auf der Erdoberfläche liegenden
geschlossenen Kurve \(s\) durch \(ds\) bezeichnet, so gilt der Satz, daß das Linienintegral der magnetischen Kraft
längs einer solchen Kurve verschwindet:
\[\int_0\mathfrak H_s\,ds=0.\]
\clearpage\SourceLeaf{67a}\HistoricalHead{62}{L}{PHYSIK. NACHLASS.}
\noindent Wählt man ein Dreieck mit den Punkten 0, 1, 2 als geschlossene Kurve \(s\), und setzt für \(\mathfrak H_s\) in den Teilintegralen
über die drei Seiten jedesmal das arithmetische Mittel aus den Werten in den Endpunkten jeder
Dreiecksseite, so ergibt sich die obige Formel unmittelbar durch Übergang zur Koordinatendarstellung.
Dasselbe Problem wird behandelt in den Artikeln 9 und 10 der \textit{Allgemeinen Theorie des Erdmagnetismus},
Resultate u.s.w. 1838, S. 11; Werke V, S. 131. Die geometrische Einkleidung fehlt dort indessen.
\par\raggedleft\textsc{Schaefer}.\par}
\clearpage\endgroup
